% Auto-extracted during the research cleanup pass. Original source is preserved unchanged.
% Source: integration.tex, analytic dense+dense integration prefix.

\section{End-to-End Integrations}
\label{sec:integrations}

This section focuses only on the dense validation pair:
a systematic dense banded outer together with the dense time-varying recursive inner.
The proof is a direct first-moment argument.
The outer contributes an expected spectrum envelope, while the inner contributes two analytic tails:
one for low/intermediate weights and one geometric tail above a logarithmic cutoff.

\subsection{Direct first-moment reduction}
\label{sec:integration-dense-dense}

Take as outer code the systematic dense log-memory construction of
Section~\ref{sec:outer-dense-logmem}, with memory
\[
M=\lceil c\log_2 n\rceil,
\]
and as inner code the dense time-varying recursive construction of
Section~\ref{sec:inner-dense}, with memory
\[
m=\lceil \gamma \log_2 n\rceil.
\]
The serial concatenation still has one random interleaver and \(O(\log n)\) taps per symbol in both stages,
hence \(O(n\log n)\) encoding work.

\begin{theorem}[Direct dense+dense first-moment bound]
\label{thm:dense-dense-direct-first-moment}
Fix \(\delta\in(0,1/2)\).
Suppose the outer ensemble satisfies
\[
\EE[\enum_{\Cone,h}] \le A_h
\qquad (1\le h\le n),
\]
and the inner satisfies
\[
p_h(\delta)
:=
\Pr[\wt(\Enc_{\mathrm{in}}(U))\le \delta n\mid \wt(U)=h]
\le B_h
\qquad (1\le h\le n),
\]
where \(U\) is uniform on the weight-\(h\) slice of \(\{0,1\}^n\).
Then
\begin{equation}
\label{eq:dense-dense-direct-first-moment}
\Pr[d_{\min}(\Csc)\le \delta n]
\;\le\;
\sum_{h=1}^{n} A_h\,B_h.
\end{equation}
\end{theorem}

\begin{proof}
By the standard first-moment reduction,
\[
\Pr[d_{\min}(\Csc)\le \delta n]
\le
\sum_{h=1}^{n}\EE[\enum_{\Cone,h}]\,p_h(\delta),
\]
because the interleaver makes every fixed weight-\(h\) outer codeword uniform on the weight-\(h\) slice before entering the inner code.
Apply the envelopes \(A_h\) and \(B_h\).
\end{proof}

\subsection{A baseline linear-distance theorem}
\label{sec:integration-dense-dense-baseline}

We now combine the global systematic-outer spectrum envelope with the two dense-inner bounds.

\begin{theorem}[Systematic dense outer + dense inner have linear distance]
\label{thm:dense-dense-integration}
Fix any \(z\in(0,\sqrt2-1)\), and define
\[
a(z):=\frac{1+z}{2}.
\]
Choose the outer memory constant \(c\) so that
\begin{equation}
\label{eq:dense-c-condition}
c>\frac{1}{\log_2(1/a(z))}.
\end{equation}
Let \(\beta_{\mathrm{out}}>1\) be the corresponding outer spectrum constant from
Corollary~\ref{lem:outer-dense-spectrum}.

Next choose constants \(\kappa>0\), \(\rho\in(0,1)\), \(\delta\in(0,1/2)\), and \(\gamma>1\) such that
\begin{equation}
\label{eq:dense-dense-parameter-conditions}
\kappa\log_2(\beta_{\mathrm{out}})<1,
\qquad
2\delta<\rho<\beta_{\mathrm{out}}^{-1},
\qquad
\gamma>1+\kappa\log_2\!\Bigl(\frac{1}{\rho}\Bigr).
\end{equation}
Then the resulting dense outer + dense inner concatenation satisfies
\[
\Pr[d_{\min}(\Csc)\le \delta n]=o(1).
\]
In particular, \(d_{\min}(\Csc)=\Omega(n)\) with high probability.
\end{theorem}

\begin{proof}
Set
\[
h_\star:=\lfloor \kappa\log_2 n\rfloor.
\]
We split the first-moment sum from Theorem~\ref{thm:dense-dense-direct-first-moment} into
\[
S_{\mathrm{low}}
:=
\sum_{h=1}^{h_\star} \EE[\enum_{\Cone,h}]\,p_h(\delta),
\qquad
S_{\mathrm{high}}
:=
\sum_{h=h_\star+1}^{n} \EE[\enum_{\Cone,h}]\,p_h(\delta).
\]

\paragraph{Outer side.}
Corollary~\ref{lem:outer-dense-spectrum} gives
\[
\EE[\enum_{\Cone,h}] \le \beta_{\mathrm{out}}^h
\qquad (1\le h\le n).
\]

\paragraph{Low-weight inner bound.}
Fix any \(\xi>0\) and let
\[
L:=\lceil(2+\xi)\delta n\rceil,
\qquad
b_{\delta,\xi}(n):=
\PP\!\left[\Binomial(L,\tfrac12)\le \delta n\right].
\]
Apply the small-weight envelope \eqref{eq:dense-inner-smallw-envelope} from
Corollary~\ref{cor:dense-inner-piecewise}:
\[
p_h(\delta)
\le
\frac{h(h-1)}{n-1}
\;+\;
\frac{\binom{L-h+1}{h}}{\binom{n-h+1}{h}}
\;+\;
h2^{-(m-1)}
\;+\;
b_{\delta,\xi}(n)
\qquad (1\le h\le h_\star).
\]
Therefore
\[
S_{\mathrm{low}}
\le
\sum_{h=1}^{h_\star}\beta_{\mathrm{out}}^h\frac{h(h-1)}{n-1}
\;+\;
\sum_{h=1}^{h_\star}\beta_{\mathrm{out}}^h
\frac{\binom{L-h+1}{h}}{\binom{n-h+1}{h}}
\;+\;
\sum_{h=1}^{h_\star}\beta_{\mathrm{out}}^h h2^{-(m-1)}
\;+\;
b_{\delta,\xi}(n)\sum_{h=1}^{h_\star}\beta_{\mathrm{out}}^h.
\]

For the first term,
\[
\sum_{h=1}^{h_\star}\beta_{\mathrm{out}}^h\frac{h(h-1)}{n-1}
\le
\frac{h_\star^2}{n}\sum_{h=1}^{h_\star}\beta_{\mathrm{out}}^h
\le
n^{\kappa\log_2(\beta_{\mathrm{out}})-1}(\log n)^{O(1)}
=o(1)
\]
by the first condition in \eqref{eq:dense-dense-parameter-conditions}.

For the fixed-tap dense outer, the tiny-weight law is more explicit than the global geometric envelope: if a message
support has band-union size \(q\) and \(\sigma\)-cluster count \(r\), then its parity contribution has law
\[
r+\Binomial(q-r,\tfrac12).
\]
In particular,
\begin{equation}
\label{eq:dense-dense-working-outer-small}
A_1^{\mathrm{out,ft}}=0
\end{equation}
deterministically.

The fixed-tap replacement exposes the same phenomenon seen in the finite-\(n\) evaluator: the late-placement term
cannot be controlled by the coarse global geometric outer envelope alone, since already \(h=1\) has inner probability
\(\approx L/n\). This low/tiny window must use the explicit dense outer low-weight law
\eqref{eq:dense-dense-working-outer-small}, rather than only \(\EE[\enum_{\Cone,h}]\le\beta_{\mathrm{out}}^h\).
With that explicit outer law, the late-placement contribution is the same summable tiny-window term used in the
preserved scalar dense diagnostics.

For the fixed-tap candidate term, since \(m=\gamma\log_2 n\),
\[
h2^{-(m-1)}\le (\log n)^{O(1)}n^{-\gamma}
\]
throughout \(h\le h_\star\), and therefore
\[
\sum_{h=1}^{h_\star}\beta_{\mathrm{out}}^h h2^{-(m-1)}
\le
(\log n)^{O(1)}n^{\kappa\log_2(\beta_{\mathrm{out}})-\gamma}
=o(1).
\]

For the binomial term, \(b_{\delta,\xi}(n)\) is exponentially small in \(n\), while
\[
\sum_{h=1}^{h_\star} \beta_{\mathrm{out}}^h
\le
(\log n)^{O(1)}\,n^{\kappa\log_2(\beta_{\mathrm{out}})},
\]
so this contribution is also \(o(1)\).
Thus
\[
S_{\mathrm{low}}=o(1).
\]

\paragraph{High-weight inner bound.}
By Corollary~\ref{cor:dense-inner-constant-geometric}, the conditions
\[
2\delta<\rho,
\qquad
\gamma>1+\kappa\log_2(1/\rho)
\]
imply that for all \(h\ge h_\star+1\),
\[
p_h(\delta)\le 3\rho^h.
\]
Therefore
\[
S_{\mathrm{high}}
\le
3\sum_{h=h_\star+1}^{n}(\beta_{\mathrm{out}}\rho)^h.
\]
Since \(\beta_{\mathrm{out}}\rho<1\), this is a geometric tail, and
\[
S_{\mathrm{high}}
\le
O(1)\,(\beta_{\mathrm{out}}\rho)^{h_\star}
=
n^{-\kappa\log_2(1/(\beta_{\mathrm{out}}\rho))+o(1)}
=o(1).
\]

Combining the low- and high-weight windows proves
\[
\Pr[d_{\min}(\Csc)\le \delta n]
\le
S_{\mathrm{low}}+S_{\mathrm{high}}
=
o(1).
\]
\end{proof}
\paragraph{What this baseline theorem buys us.}
Theorem~\ref{thm:dense-dense-integration} is fully analytic and already matches the qualitative picture from the
enumerator experiments:
\begin{itemize}
\item the systematic outer removes collapse to tiny intermediate weights,
\item the outer contributes a global geometric spectrum envelope,
\item the inner handles small weights by the run-aware bound and large weights by a geometric cutoff.
\end{itemize}
This is the baseline dense+dense theorem. The next quantitative question is how close one can push \(\delta\) toward
the random-linear benchmark while remaining within the admissible parameter region
\eqref{eq:dense-dense-parameter-conditions}.

\begin{corollary}[Explicit admissible dense+dense parameter region]
\label{cor:dense-dense-explicit-region}
Fix any \(z\in(0,\sqrt2-1)\), and choose
\[
c>\frac{1}{\log_2(2/(1+z))}.
\]
Let \(\kappa>0\), \(\rho\in(0,z)\), \(\delta\in(0,\rho/2)\), and \(\gamma>1\) satisfy
\[
\kappa\log_2(1/z)<1,
\qquad
\gamma>1+\kappa\log_2\!\Bigl(\frac{1}{\rho}\Bigr).
\]
Then the systematic dense outer with memory \(M=\lceil c\log_2 n\rceil\) concatenated with the dense inner of memory
\(m=\lceil \gamma\log_2 n\rceil\) satisfies
\[
\Pr[d_{\min}(\Csc)\le \delta n]=o(1).
\]
\end{corollary}

\begin{proof}
For all sufficiently large \(n\), Corollary~\ref{lem:outer-dense-spectrum} gives
\[
\EE[\enum_{\Cone,h}] \le z^{-h},
\]
so Theorem~\ref{thm:dense-dense-integration} applies with \(\beta_{\mathrm{out}}=z^{-1}\).
The stated conditions are exactly the theorem's hypotheses rewritten in terms of \(z\).
\end{proof}

\begin{corollary}[One concrete baseline choice]
\label{cor:dense-dense-concrete-choice}
Choose
\[
z=\frac13,
\qquad
\kappa=\frac12,
\qquad
\rho=\frac14.
\]
If
\[
c>\frac{1}{\log_2(3/2)},
\qquad
\gamma>2,
\qquad
\delta<\frac18,
\]
then the systematic dense outer with memory \(M=\lceil c\log_2 n\rceil\) and the dense inner with memory
\(m=\lceil \gamma\log_2 n\rceil\) satisfy
\[
\Pr[d_{\min}(\Csc)\le \delta n]=o(1).
\]
\end{corollary}

\begin{proof}
For \(z=\tfrac13\), we have \(a(z)=\tfrac23\), so the outer condition becomes
\[
c>\frac{1}{\log_2(3/2)}.
\]
Also,
\[
\kappa\log_2(1/z)=\frac12\log_2 3<1,
\]
and
\[
1+\kappa\log_2(1/\rho)
=
1+\frac12\log_2 4
=2.
\]
Finally, \(\delta<1/8\) implies \(2\delta<1/4=\rho<z=1/3\).
Thus Corollary~\ref{cor:dense-dense-explicit-region} applies.
\end{proof}

\begin{remark}[Memory-vs.-blocklength interpretation]
\label{rem:dense-dense-memory-scaling}
The concrete choice in Corollary~\ref{cor:dense-dense-concrete-choice} can be restated directly in terms of the
log-memory parameter.
If the outer memory is
\[
M=\lceil c\log_2 n\rceil
\]
with
\[
c>\frac{1}{\log_2(3/2)}\approx 1.71,
\]
and the inner memory is
\[
m=\lceil \gamma\log_2 n\rceil
\qquad\text{with}\qquad \gamma>2,
\]
then the dense serial concatenation has linear minimum distance with high probability for every \(\delta<1/8\).
Equivalently, fixing a memory budget \(M\), the baseline proof certifies linear distance throughout the regime
\[
n \le 2^{M/c}
\]
for any fixed \(c>1/\log_2(3/2)\).
This does not yet match the stronger finite-length trend suggested by the enumerator, but it gives a rigorous
logarithmic-memory scaling window from which sharper dense bounds can now be pursued.
\end{remark}

\subsection{A linear-weight exponent criterion}
\label{sec:integration-dense-dense-linear}

The next sharpening step is to replace the coarse high-weight geometric tail by a direct comparison of the outer and
inner linear-weight exponents.

\begin{theorem}[Dense+dense linear-weight criterion]
\label{thm:dense-dense-linear-weight-criterion}
Fix constants \(\eta_0\in(0,1)\), \(\delta\in(0,1/2)\), \(\xi>0\), and \(\theta\in(0,1-\eta_0)\).
For \(\eta\in[\eta_0,1-\theta]\), let
\[
\Phi_{\mathrm{out}}^{\mathrm{sys}}(\eta)
\]
be the outer exponent from Theorem~\ref{thm:outer-dense-linear-exponent}, and let
\[
E_{\mathrm{in}}(\eta;\theta,\delta,\xi)
\]
be the dense-inner exponent from Theorem~\ref{prop:dense-inner-linear-exponent}.
Assume there exists \(\tau>0\) such that
\begin{equation}
\label{eq:dense-dense-linear-gap}
\Phi_{\mathrm{out}}^{\mathrm{sys}}(\eta)
\;-\;
E_{\mathrm{in}}(\eta;\theta,\delta,\xi)
\;\le\;
-\tau
\qquad\text{for every }\eta\in[\eta_0,1-\theta].
\end{equation}
Then the total first-moment contribution from outer weights \(h\ge \eta_0 n\) satisfies
\[
\sum_{h=\lceil \eta_0 n\rceil}^{\lfloor (1-\theta)n\rfloor}
\EE[\enum_{\Cone,h}]\,p_h(\delta)
\;\le\;
n^{O(1)}\,2^{-\tau n}.
\]
\end{theorem}

\begin{proof}
For \(h=\lfloor \eta n\rfloor\) with \(\eta\in[\eta_0,1-\theta]\), Theorem~\ref{thm:outer-dense-linear-exponent} gives
\[
\EE[\enum_{\Cone,h}]
\le
n^{O(1)}\,2^{n\Phi_{\mathrm{out}}^{\mathrm{sys}}(\eta)}.
\]
Theorem~\ref{prop:dense-inner-linear-exponent} gives
\[
p_h(\delta)
\le
n^{O(1)}\,2^{-nE_{\mathrm{in}}(\eta;\theta,\delta,\xi)}.
\]
Multiplying and using \eqref{eq:dense-dense-linear-gap},
\[
\EE[\enum_{\Cone,h}]\,p_h(\delta)
\le
n^{O(1)}\,2^{-\tau n}.
\]
There are only \(O(n)\) such weights \(h\), so summing over
\(h\in[\eta_0 n,(1-\theta)n]\) preserves the same exponential rate.
\end{proof}

\paragraph{Meaning of the criterion.}
Theorem~\ref{thm:dense-dense-linear-weight-criterion} isolates the real quantitative problem left by the baseline proof.
Once the low and intermediate weights are handled by the coarse argument, the only remaining analytic task is to show
that the outer exponent \(\Phi_{\mathrm{out}}^{\mathrm{sys}}\) sits below the dense-inner exponent
\(E_{\mathrm{in}}\) on the linear-weight window of interest.

\begin{theorem}[A concrete moderate-to-high-weight exponent gap]
\label{prop:dense-dense-highweight-gap}
Take
\[
\delta=0.109,
\qquad
\theta=0.005,
\qquad
\xi=8.
\]
Then direct evaluation of the explicit rate functions from
Corollary~\ref{cor:outer-dense-highweight-exponent} and
Theorem~\ref{prop:dense-inner-linear-exponent} shows that
\[
\Phi_{\mathrm{out}}^{\mathrm{sys}}(\eta)
\;
-\;
E_{\mathrm{in}}(\eta;0.005,0.109,8)
\le -0.003
\qquad
\text{for every }\eta\in[\eta_{\mathrm{crit}},0.99],
\]
where
\[
\eta_{\mathrm{crit}}:=1-\frac{1}{\sqrt2}.
\]
Consequently,
\[
\sum_{h=\lceil \eta_{\mathrm{crit}}n\rceil}^{\lfloor 0.99n\rfloor}
\EE[\enum_{\Cone,h}]\,p_h(0.109)
\le
n^{O(1)}\,2^{-0.003n}.
\]
\end{theorem}

\begin{proof}
The inner exponent is
\[
E_{\mathrm{in}}(\eta;0.005,0.109,8)
=
\min\!\left\{
\Psi_{\mathrm{run}}^{(w)}(\eta,0.005),\
1-H_2(0.109)
\right\}.
\]
The explicit binomial term equals
\[
1-H_2(0.109)\approx 0.5031.
\]
On \([\eta_{\mathrm{crit}},0.99]\), Corollary~\ref{cor:outer-dense-highweight-exponent} gives
\[
\Phi_{\mathrm{out}}^{\mathrm{sys}}(\eta)=H_2(\eta)-\frac12.
\]
A direct numerical evaluation of the explicit formulas on \([\eta_{\mathrm{crit}},0.99]\) then shows that the gap
\[
\Phi_{\mathrm{out}}^{\mathrm{sys}}(\eta)-E_{\mathrm{in}}(\eta;0.005,0.109,8)
\]
never exceeds \(-0.003\); the worst point occurs near \(\eta\approx 0.5\).
This verification is reproduced by the companion script
\texttt{scripts/verify\_dense\_claims.py}, which for the present parameters reports a worst sampled gap
of approximately \(-0.0031\) near \(\eta=0.5000\).
The final summation bound is then exactly Theorem~\ref{thm:dense-dense-linear-weight-criterion}.
\end{proof}

\begin{theorem}[A concrete 10.9-percent dense+dense theorem]
\label{thm:dense-dense-explicit-constant}
Assume
\[
c>\frac{1}{\log_2(3/2)},
\qquad
\gamma>2.
\]
Then the systematic dense outer with memory
\[
M=\lceil c\log_2 n\rceil
\]
concatenated with the dense inner of memory
\[
m=\lceil \gamma\log_2 n\rceil
\]
satisfies
\[
\Pr[d_{\min}(\Csc)\le 0.109 n]=o(1).
\]
\end{theorem}

\begin{proof}
Set
\[
\delta=0.109,
\qquad
\kappa=\frac12,
\qquad
\rho=\frac14,
\qquad
\theta=0.005,
\qquad
\xi=8,
\qquad
h_0:=\lfloor \kappa\log_2 n\rfloor,
\]
and let \(\eta_{\mathrm{crit}}:=1-2^{-1/2}\).
Split the first-moment sum into four windows:
\[
S_{\mathrm{tiny}}
:=
\sum_{h=1}^{h_0}\EE[\enum_{\Cone,h}]\,p_h(0.109),
\]
\[
S_{\mathrm{low}}
:=
\sum_{h=h_0+1}^{\lfloor \eta_{\mathrm{crit}}n\rfloor}\EE[\enum_{\Cone,h}]\,p_h(0.109),
\]
\[
S_{\mathrm{lin}}
:=
\sum_{h=\lceil \eta_{\mathrm{crit}}n\rceil}^{\lfloor 0.99n\rfloor}\EE[\enum_{\Cone,h}]\,p_h(0.109),
\]
and
\[
S_{\mathrm{top}}
:=
\sum_{h=\lceil 0.99n\rceil}^{n}\EE[\enum_{\Cone,h}]\,p_h(0.109).
\]

For \(S_{\mathrm{tiny}}\), the proof of Corollary~\ref{cor:dense-dense-concrete-choice} already applies verbatim,
since \(0.109<1/8\), and gives \(S_{\mathrm{tiny}}=o(1)\).

For \(S_{\mathrm{low}}\), Corollary~\ref{cor:outer-dense-lowlinear-envelope} gives
\[
\EE[\enum_{\Cone,h}] \le n^{O(1)}(1+\sqrt2)^h
\qquad(h\le \eta_{\mathrm{crit}}n).
\]
Also, Corollary~\ref{cor:dense-inner-constant-geometric} applies with \(\rho=1/4\), \(\kappa=1/2\), and \(\gamma>2\),
so for all \(h\ge h_0+1\),
\[
p_h(0.109)\le 3\cdot 4^{-h}.
\]
Therefore
\[
S_{\mathrm{low}}
\le
n^{O(1)}\sum_{h=h_0+1}^{\lfloor \eta_{\mathrm{crit}}n\rfloor}
\left(\frac{1+\sqrt2}{4}\right)^h.
\]
Since \((1+\sqrt2)/4<1\), this is a geometric tail beginning at \(h_0=\tfrac12\log_2 n+O(1)\), hence
\[
S_{\mathrm{low}}
\le
n^{-\frac12\log_2(4/(1+\sqrt2))+o(1)}
=
o(1).
\]

For \(S_{\mathrm{lin}}\), Theorem~\ref{prop:dense-dense-highweight-gap} gives
\[
S_{\mathrm{lin}}\le n^{O(1)}2^{-0.003n}=o(1).
\]

For \(S_{\mathrm{top}}\), use only the outer exponent and the trivial bound \(p_h(0.109)\le 1\).
If \(h=\lfloor \eta n\rfloor\) with \(\eta\in[0.99,1]\), then Corollary~\ref{cor:outer-dense-highweight-exponent} gives
\[
\Phi_{\mathrm{out}}^{\mathrm{sys}}(\eta)=H_2(\eta)-\frac12
\le
H_2(0.99)-\frac12
< -0.41.
\]
Hence Theorem~\ref{thm:outer-dense-linear-exponent} gives
\[
\EE[\enum_{\Cone,h}]
\le
n^{O(1)}2^{-0.41n},
\]
uniformly over \(h\in[0.99n,n]\), and so
\[
S_{\mathrm{top}}
\le
n^{O(1)}2^{-0.41n}
=
o(1).
\]

Combining the four windows proves
\[
\Pr[d_{\min}(\Csc)\le 0.109 n]
\le
S_{\mathrm{tiny}}+S_{\mathrm{low}}+S_{\mathrm{lin}}+S_{\mathrm{top}}
=
o(1).
\]
\end{proof}
